% from ega4-8-fr.tex
\oldpage[IV-3]{10}
de $T$ (resp. $T'$). Soit $(S_\lambda,u_{\lambda\mu})$ un système projectif
(non nécessairement filtrant) dans $\mathbf C_T$, et posons
$S'_\lambda=S_\lambda\times_T T'$,
$u'_{\lambda\mu}=u_{\lambda\mu}\times 1_{T'}$, de sorte que
$(S'_\lambda,u'_{\lambda\mu})$ est un système projectif dans
$\mathbf C_{T'}$. Alors, si $(S,u_\lambda)$ est une limite projective de
$(S_\lambda,u_{\lambda\mu})$ dans $\mathbf C_T$,
$(S\times_T T',u_\lambda\times 1_{T'})$ est une limite projective de
$(S'_\lambda,u'_{\lambda\mu})$ dans $\mathbf C_{T'}$.
\end{lemma}

On a par hypothèse, pour tout $\lambda$, un diagramme commutatif
\[
\xymatrix{
  S' \ar[r]^{u'_\lambda} \ar[d]_p
    & S'_\lambda \ar[r]^{h'_\lambda} \ar[d]_{p_\lambda}
    & T' \ar[d]^q \\
  S \ar[r]_{u_\lambda} & S_\lambda \ar[r]_{h_\lambda} & T
}
\]
où l'on a posé $S'=S\times_T T'$, $u'_\lambda=u_\lambda\times 1_{T'}$,
$h'_\lambda=h_\lambda\times 1_{T'}$. Soient $Y$ un $T'$-objet,
$g':Y\to T'$ le morphisme correspondant, et considérons un système
projectif de $T'$-morphismes $w'_\lambda:Y\to S'_\lambda$. Alors $Y$ est un
$T$-objet pour le morphisme $g=q\circ g'$, et les
$w_\lambda=p_\lambda\circ w'_\lambda$ sont des $T$-morphismes puisque
$h_\lambda\circ w_\lambda=h_\lambda\circ p_\lambda\circ w'_\lambda
=q\circ h'_\lambda\circ w'_\lambda=q\circ g'$ par hypothèse. En outre, on
vérifie aussitôt que $(w_\lambda)$ est un système projectif. Il existe donc
par hypothèse un $T$-morphisme $w:Y\to S$ et un seul tel que
$u_\lambda\circ w=w_\lambda$ pour tout $\lambda$. Par définition du produit
fibré, il y a un unique $T'$-morphisme $w':Y\to S'$ tel que
$p\circ w'=w$. On a alors
$u_\lambda\circ p\circ w'=u_\lambda\circ w=w_\lambda
=p_\lambda\circ w'_\lambda$, ce qui s'écrit aussi
$p_\lambda\circ u'_\lambda\circ w'=p_\lambda\circ w'_\lambda$~; d'autre
part, en écrivant que $w'$ est un $T'$-morphisme, il vient
$h'_\lambda\circ u'_\lambda\circ w'=g'=h'_\lambda\circ w'_\lambda$. La
définition de $S'_\lambda$ comme produit fibré $S_\lambda\times_T T'$ donne
donc $u'_\lambda\circ w'=w'_\lambda$, et il est immédiat que $w'$ est
l'unique $T'$-morphisme vérifiant ces relations, d'où le lemme.

\begin{remark}[8.2.7]
\label{IV.8.2.7-fr}
Étant donné un espace annelé quelconque $S$, les limites inductives par
rapport à un ensemble préordonné \emph{quelconque} $L$ (non nécessairement
filtrant) existent dans la catégorie des $\mathcal O_S$-Algèbres
commutatives, puisque la limite inductive filtrante existe d'après (8.2.1) et
d'autre part, pour deux homomorphismes de $\mathcal O_S$-Algèbres
$\mathscr A\to\mathscr B$, $\mathscr A\to\mathscr C$, le produit tensoriel
$\mathscr B\otimes_{\mathscr A}\mathscr C$ est la «~somme amalgamée~»
correspondante dans cette catégorie.

Lorsque $S$ est un préschéma, on sait que le produit tensoriel
$\mathscr B\otimes_{\mathscr A}\mathscr C$ est une $\mathcal O_S$-Algèbre
quasi-cohérente lorsqu'il en est ainsi de $\mathscr A$, $\mathscr B$ et
$\mathscr C$ (I, 1.3.13)~; on en conclut que, dans la catégorie des
$\mathcal O_S$-Algèbres \emph{quasi-cohérentes}, les limites inductives pour
un ensemble préordonné quelconque d'indices existent toujours. Cela permet de
généraliser la définition d'une limite projective de préschémas et les prop.
(8.2.3) et (8.2.5) au cas où l'ensemble préordonné $L$ n'est pas
nécessairement filtrant.
\end{remark}

\begin{env}[8.2.8]
\label{IV.8.2.8-fr}
Les notations étant celles de (8.2.2), posons
$u_{\lambda\mu}=(\psi_{\lambda\mu},\theta_{\lambda\mu})$ et
$u_\lambda=(\psi_\lambda,\theta_\lambda)$~; $(S_\lambda,\psi_{\lambda\mu})$
est donc un système projectif d'espaces topologiques et $(\psi_\lambda)$ un
système projectif d'applications continues $S\to S_\lambda$, des espaces
sous-jacents respectivement aux préschémas $S$ et $S_\lambda$.
\end{env}

\begin{proposition}[8.2.9]
\label{IV.8.2.9-fr}
Avec les notations de (8.2.8), la limite projective du système projectif
$(\psi_\lambda)$ d'applications continues est un homéomorphisme de l'espace
sous-jacent à $S$ sur la limite projective du système projectif
$(S_\lambda,\psi_{\lambda\mu})$ d'espaces topologiques.
\end{proposition}

